\section{Introduction}

Large Language Model (LLM) judges have become central infrastructure for NLP evaluation and preference-based alignment training.
Systems such as AlpacaEval, MT-Bench, and Chatbot Arena rely on LLM judges to score and rank model outputs, and the resulting leaderboards guide both academic research and industry deployment.
Reinforcement Learning from Human Feedback (RLHF) and Direct Preference Optimization (DPO) pipelines consume judge-generated preferences as reward signals that shape model behavior during alignment.
The reliability of these judgments carries direct consequences: it determines which models the community considers superior and which response patterns are reinforced.
As automated judges increasingly mediate both evaluation and training, any systematic bias in their assessments propagates through the entire model development pipeline.

Recent work has converged on a striking finding: \emph{style bias} is the dominant source of unreliability in LLM judges, far outweighing other documented biases such as position or verbosity preference.
\citet{murugadoss2025sos} show that style features alone explain judge scores with $R = 0.999$, far exceeding the explanatory power of content features.
\citet{li2026preference} demonstrate that 76--92\% of observed preference signals originate from surface-level style cues rather than substantive quality differences.
\citet{soumik2026judging} report bias scores of 0.76--0.92 across five different judges.
Despite this consensus on severity, existing approaches share four critical gaps: (1)~they treat style as a single dimension, typically reducing it to response length; (2)~they rely on correlational analyses that cannot establish causal direction; (3)~population-level corrections such as length-controlled AlpacaEval \citep{dubois2025lc} cannot diagnose per-instance bias; and (4)~no existing method provides black-box, post-hoc, per-axis causal detection.
Recent multi-feature frameworks such as CASSE \citep{casse2026concise} and FairJudge \citep{wang2026fairjudge} begin to decompose style into finer dimensions, but neither provides causal identification of per-axis effects.
A further methodological concern compounds these limitations: standard Likert pointwise scoring (integer scales from 1 to 10) exhibits a ceiling effect that masks real bias.
On our data, 64\% of counterfactual pairs receive identical Likert scores ($\Delta = 0$), and this integer ceiling effect leaves Likert insensitive to style interventions that pairwise comparison readily detects.

We propose Causal Style Decomposition (CSD), a framework that addresses all four gaps.
Drawing on causal inference methods for text \citep{veitch2020adapting}, CSD operates in three stages: binary style counterfactual generation along a single axis (e.g., formal vs.\ casual), content preservation verification through Natural Language Inference (NLI) entailment and BERTScore filtering, and causal bias measurement via position-controlled pairwise comparison.
Because CSD requires no access to judge internals and applies post-hoc to any black-box judge, it complements existing population-level approaches and provides finer-grained, per-instance diagnostics.
Our experiments reveal what we term the \emph{Likert Blindspot}: on the same formality data, Likert scoring yields a corrected Average Treatment Effect (ATE) of $+0.102$ (Cohen's $d = 0.104$, $p = 0.332$, $n = 88$), a non-significant result, while pairwise comparison shows a formal win rate of 69.1\% (Wilson CI: $[55.97\%, 79.72\%]$, sign test $p = 0.006$, $n = 55$ consistent pairs).
Figure~\ref{fig:likert_blindspot} illustrates this divergence.
Unlike LC-AlpacaEval \citep{dubois2025lc}, which applies a single population-level length regression, CSD decomposes bias along multiple axes with per-instance causal intervention.

Our contributions are as follows:
\begin{itemize}
    \item \textbf{A black-box causal framework} for detecting per-axis style bias in LLM judges through binary counterfactual generation and position-controlled pairwise comparison (Section~\ref{sec:method}).
    \item \textbf{The Likert Blindspot phenomenon:} empirical evidence that standard pointwise scoring systematically fails to detect style biases that pairwise comparison reveals, with formality bias undetectable by Likert ($p = 0.332$) but significant under pairwise measurement (69.1\%, $p = 0.006$) (Section~\ref{sec:experiments}, Figure~\ref{fig:likert_blindspot}).
    \item \textbf{Evidence of axis-dependent bias:} \textit{[to be updated with full experimental results across multiple judges and style axes]}.
    \item \textbf{Practical implications for evaluation methodology:} style bias audits should employ pairwise comparison rather than Likert scoring, and bias must be assessed per axis rather than as a single ``style'' factor (Section~\ref{sec:discussion}).
\end{itemize}
